\section{Related work and a practical map}
\label{sec:related}

Semantic agreement between informal mathematics and formal statements is now an
explicit target across several lines of work.  Statement-level methods include
learned or LLM-based critics such as FormalAlign, CriticLean, and Beyond
Compilation \citep{lu2025formalalign,peng2025criticlean,zhang2026beyondcompilation};
reference-based equivalence checks such as BEq/BEq+
\citep{liu2025beq,poiroux2025reliable}; and compiler-grounded tests such as
Theorem Testing and ShadowBench/SA-Pass
\citep{kim2026theoremtesting,han2026shadowbench}.  These methods help test
whether a generated statement agrees with a target when an appropriate
reference, critic, dependent theorem set, or family of shadow statements is
available.

Project-scale systems address additional questions once the source and the Lean
development both evolve.  EconCSLib provides paper-facing Lean interfaces,
dependency graphs, validation reports, and a human review dashboard
\citep{garg2026econcs}.  Lean Atlas/Lean Compass identifies project declarations
whose semantics can affect selected target theorems and presents them in an
interactive review environment \citep{yanahama2026leanatlas}.  ATLAS/AutoformBot
supports comparison of textbook statements with Lean formalizations, dependency
inspection, and automated faithfulness assessments \citep{rammal2026scale}.
LeanMarathon, FormaTheoria, and LeanArchitect address target review,
long-horizon literature formalization, provenance, and synchronization between
informal and formal project descriptions
\citep{zhang2026leanmarathon,nie2026formatheoria,zhu2026leanarchitect}.

Agent-facing Lean infrastructure sits underneath these review systems.  LeanDojo
provides programmatic interaction and retrieval infrastructure for theorem
proving agents \citep{yang2023leandojo}.  Our open-source \texttt{leanq} tool
queries the elaborated Lean environment for declarations, theorem types, axioms,
dependencies, reverse dependencies, and graph slices
\citep{aiqlean2026tools}.  In our workflow, these tools reduce the amount of Lean
state an LLM or human must recover from source text alone; they do not decide
whether the resulting theorem matches a paper.

\begin{table}[t]
\centering
\small
\caption{A practical map for newcomers.  The systems overlap; each row names a
question for which the cited work provides a useful starting point.}
\label{tab:tool-map}
\begin{tabularx}{\linewidth}{@{}P{0.25\linewidth}P{0.35\linewidth}Y@{}}
\toprule
Question & Representative work & Main aid \\
\midrule
Does this statement match the intended theorem? & BEq/BEq+, CriticLean, ShadowBench & equivalence tests, semantic critics, shadow statements \\
What should a human inspect? & Lean Atlas / Compass & review cone and interactive dependency view \\
How can an author review a paper translation? & EconCSLib & paper-to-Lean dashboard and human judgments \\
How do agents query Lean? & LeanDojo, \texttt{leanq} & interaction, retrieval, elaborated statements, dependencies \\
How should review fit into a long project? & LeanMarathon, FormaTheoria, LeanArchitect & target review, provenance, reconciliation, blueprint maintenance \\
\bottomrule
\end{tabularx}
\end{table}

The local dashboard in this paper grew from one Davis--Kahan formalization and
overlaps with several of these systems.  We use the comparison to organize
practical lessons and to point newcomers toward existing tools rather than to
make a priority claim for semantic review or review dashboards.
